% GENERATED FILE. Edit canonical lessons, then run npm run generate:books.
% canonical-source-hash: fnv1a64:cbeb5c5c94078671
% canonical-lessons: MW-C225-gathri, MW-C225-jholi, MW-C225-masak, MW-C225-bistar, MW-C225-kamli

\chapter{Packing for the Road}
\label{ch:mw-packing-for-the-road}
\hlchaptermodality{225}

\begin{chapteropening}
\textbf{By the end of this chapter:} \emph{I can name what I will pack for a journey.}

Complete the last lesson of chapter 225: I can name what I will pack for a journey.
\end{chapteropening}

\section[\texorpdfstring{\mw{गठड़ी} (gaṭhṛī)}{गठड़ी (gaṭhṛī)}]{\mw{गठड़ी} (gaṭhṛī) — a bundle (tied in cloth)}

\label{lesson:MW-C225-gathri}



\begin{quote}
Before the new one: say the Marwadi for to leap up, then the Marwadi for a kettledrum.
\end{quote}

\subsection*{You'll want to know: \mw{गठड़ी}}
\textbf{\mw{गठड़ी}} — \emph{gaṭhṛī} — \textquotedblleft{}a bundle (tied in cloth)\textquotedblright{}.

\begin{quote}
\mw{म्हारी} \mw{गठड़ी} \mw{भारी} \mw{है।}
\end{quote}

\begin{quote}
\emph{mhārī gaṭhṛī bhārī hai.}
\end{quote}

\begin{quote}
My bundle is heavy.
\end{quote}

A feminine noun.

\begin{cousinweb}
\textbf{\mw{गठड़ी}} (\emph{gaṭhṛī}) is inherited from Sanskrit \emph{grathna}, \textquotedblleft{}a bunch\textquotedblright{}, from the root \emph{granth}, \textquotedblleft{}to tie\textquotedblright{}; CDIAL lists Hindi \emph{gaṭhṛī}, \textquotedblleft{}a bundle\textquotedblright{}. CDIAL gives no Marwari form in this entry (CDIAL 4348).
\end{cousinweb}

\begin{grammarlens}[title={What you've built}]
One more word for saying what comes next.
\end{grammarlens}

\subsection*{Your turn}
\begin{itemize}
\raggedright
  \item \emph{Say it:} \emph{gaṭhṛī}
  \item \emph{Say it:} \emph{gaṭhṛī}, once more
  \item \emph{Say it:} \emph{mhārī gaṭhṛī bhārī hai.} — My bundle is heavy.
  \item \emph{Recall:} say the Marwadi for to leap up, then the Marwadi for a kettledrum, then say \emph{gaṭhṛī} again
\end{itemize}

\begin{tcolorbox}[breakable,colback=teal!4,colframe=teal!35!black,title={Before you move on}]
What does \mw{गठड़ी} mean? (\textquotedblleft{}A bundle (tied in cloth)\textquotedblright{}.) Say it once more.
\end{tcolorbox}
\section[\texorpdfstring{\mw{झोळी} (jhoḷī)}{झोळी (jhoḷī)}]{\mw{झोळी} (jhoḷī) — a cloth bag (carried on the shoulder)}

\label{lesson:MW-C225-jholi}



\begin{quote}
Before the new one: say the Marwadi for a kettledrum, then the Marwadi for a bundle.
\end{quote}

\subsection*{You'll want to know: \mw{झोळी}}
\textbf{\mw{झोळी}} — \emph{jhoḷī} — \textquotedblleft{}a cloth bag (carried on the shoulder)\textquotedblright{}.

\begin{quote}
\mw{हूं} \mw{काल} \mw{झोळी} \mw{मांय} \mw{रोटी} \mw{राखसूं।}
\end{quote}

\begin{quote}
\emph{hū̃ kāl jhoḷī mā̃y roṭī rākhsū̃.}
\end{quote}

\begin{quote}
Tomorrow I'll put flatbread in my bag.
\end{quote}

A feminine noun.

\begin{cousinweb}
\textbf{\mw{झोळी}} (\emph{jhoḷī}) is inherited from a reconstructed \emph{*jhōla}, \textquotedblleft{}a bag\textquotedblright{}; CDIAL lists Hindi \emph{jholī} and Gujarati \emph{jhoḷī}, \textquotedblleft{}a bag\textquotedblright{}. CDIAL gives no Marwari form in this entry (CDIAL 5415).
\end{cousinweb}

\begin{grammarlens}[title={What you've built}]
One more word for saying what comes next.
\end{grammarlens}

\subsection*{Your turn}
\begin{itemize}
\raggedright
  \item \emph{Say it:} \emph{jhoḷī}
  \item \emph{Say it:} \emph{jhoḷī}, once more
  \item \emph{Say it:} \emph{hū̃ kāl jhoḷī mā̃y roṭī rākhsū̃.} — Tomorrow I'll put flatbread in my bag.
  \item \emph{Recall:} say the Marwadi for a kettledrum, then the Marwadi for a bundle, then say \emph{jhoḷī} again
\end{itemize}

\begin{tcolorbox}[breakable,colback=teal!4,colframe=teal!35!black,title={Before you move on}]
What does \mw{झोळी} mean? (\textquotedblleft{}A cloth bag (carried on the shoulder)\textquotedblright{}.) Say it once more.
\end{tcolorbox}
\section[\texorpdfstring{\mw{मसक} (masak)}{मसक (masak)}]{\mw{मसक} (masak) — a water-skin}

\label{lesson:MW-C225-masak}



\begin{quote}
Before the new one: say the Marwadi for a bundle, then the Marwadi for a cloth bag.
\end{quote}

\subsection*{You'll want to know: \mw{मसक}}
\textbf{\mw{मसक}} — \emph{masak} — \textquotedblleft{}a water-skin\textquotedblright{}.

\begin{quote}
\mw{हूं} \mw{मसक} \mw{मांय} \mw{पाणी} \mw{भरसूं।}
\end{quote}

\begin{quote}
\emph{hū̃ masak mā̃y pāṇī bharsū̃.}
\end{quote}

\begin{quote}
I'll fill the water-skin.
\end{quote}

A feminine noun.

\begin{cousinweb}
\textbf{\mw{मसक}} (\emph{masak}) is usually traced to Persian \emph{mašk}, \textquotedblleft{}a leather bag for water\textquotedblright{}; Marwadi breaks up its final \emph{-šk} with a vowel.
\end{cousinweb}

\begin{grammarlens}[title={What you've built}]
One more word for saying what comes next.
\end{grammarlens}

\subsection*{Your turn}
\begin{itemize}
\raggedright
  \item \emph{Say it:} \emph{masak}
  \item \emph{Say it:} \emph{masak}, once more
  \item \emph{Say it:} \emph{hū̃ masak mā̃y pāṇī bharsū̃.} — I'll fill the water-skin.
  \item \emph{Recall:} say the Marwadi for a bundle, then the Marwadi for a cloth bag, then say \emph{masak} again
\end{itemize}

\begin{tcolorbox}[breakable,colback=teal!4,colframe=teal!35!black,title={Before you move on}]
What does \mw{मसक} mean? (\textquotedblleft{}A water-skin\textquotedblright{}.) Say it once more.
\end{tcolorbox}
\section[\texorpdfstring{\mw{बिस्तर} (bistar)}{बिस्तर (bistar)}]{\mw{बिस्तर} (bistar) — bedding; a bedroll}

\label{lesson:MW-C225-bistar}



\begin{quote}
Before the new one: say the Marwadi for a cloth bag, then the Marwadi for a water-skin.
\end{quote}

\subsection*{You'll want to know: \mw{बिस्तर}}
\textbf{\mw{बिस्तर}} — \emph{bistar} — \textquotedblleft{}bedding; a bedroll\textquotedblright{}.

\begin{quote}
\mw{हूं} \mw{काल} \mw{बिस्तर} \mw{लेसूं।}
\end{quote}

\begin{quote}
\emph{hū̃ kāl bistar lesū̃.}
\end{quote}

\begin{quote}
Tomorrow I'll take my bedding.
\end{quote}

A masculine noun.

\begin{cousinweb}
\textbf{\mw{बिस्तर}} (\emph{bistar}) is usually traced to Persian \emph{bistar}, \textquotedblleft{}a bed, bedding\textquotedblright{}.
\end{cousinweb}

\begin{grammarlens}[title={What you've built}]
One more word for saying what comes next.
\end{grammarlens}

\subsection*{Your turn}
\begin{itemize}
\raggedright
  \item \emph{Say it:} \emph{bistar}
  \item \emph{Say it:} \emph{bistar}, once more
  \item \emph{Say it:} \emph{hū̃ kāl bistar lesū̃.} — Tomorrow I'll take my bedding.
  \item \emph{Recall:} say the Marwadi for a cloth bag, then the Marwadi for a water-skin, then say \emph{bistar} again
\end{itemize}

\begin{tcolorbox}[breakable,colback=teal!4,colframe=teal!35!black,title={Before you move on}]
What does \mw{बिस्तर} mean? (\textquotedblleft{}Bedding; a bedroll\textquotedblright{}.) Say it once more.
\end{tcolorbox}
\section[\texorpdfstring{\mw{कामळी} (kāmḷī)}{कामळी (kāmḷī)}]{\mw{कामळी} (kāmḷī) — a woollen blanket}

\label{lesson:MW-C225-kamli}



\begin{quote}
Before the new one: say the Marwadi for a water-skin, then the Marwadi for bedding.
\end{quote}

\subsection*{You'll want to know: \mw{कामळी}}
\textbf{\mw{कामळी}} — \emph{kāmḷī} — \textquotedblleft{}a woollen blanket\textquotedblright{}.

\begin{quote}
\mw{रात} \mw{ठंडी} \mw{है}, \mw{हूं} \mw{कामळी} \mw{लेसूं।}
\end{quote}

\begin{quote}
\emph{rāt ṭhaṇḍī hai, hū̃ kāmḷī lesū̃.}
\end{quote}

\begin{quote}
The night is cold, so I'll take a blanket.
\end{quote}

A feminine noun.

\begin{cousinweb}
\textbf{\mw{कामळी}} (\emph{kāmḷī}) is inherited from Sanskrit \emph{kambala}, \textquotedblleft{}a woollen blanket\textquotedblright{}; CDIAL lists Hindi \emph{kamlī}, \textquotedblleft{}a small blanket\textquotedblright{}, and Gujarati \emph{kāmḷī}, \textquotedblleft{}a blanket\textquotedblright{}. CDIAL gives no Marwari form in this entry (CDIAL 2771).
\end{cousinweb}

\begin{grammarlens}[title={What you've built}]
One more word for saying what comes next.
\end{grammarlens}

\subsection*{Your turn}
\begin{itemize}
\raggedright
  \item \emph{Say it:} \emph{kāmḷī}
  \item \emph{Say it:} \emph{kāmḷī}, once more
  \item \emph{Say it:} \emph{rāt ṭhaṇḍī hai, hū̃ kāmḷī lesū̃.} — The night is cold, so I'll take a blanket.
  \item \emph{Recall:} say the Marwadi for a water-skin, then the Marwadi for bedding, then say \emph{kāmḷī} again
\end{itemize}

\begin{tcolorbox}[breakable,colback=teal!4,colframe=teal!35!black,title={Before you move on}]
What does \mw{कामळी} mean? (\textquotedblleft{}A woollen blanket\textquotedblright{}.) Say it once more.
\end{tcolorbox}
